%------------------------------------------------------------------------------%
\documentclass[fleqn,utf8,aspectratio=169,12pt]{beamer}

\usepackage{gaal}

\title{Diagonalização de Matrizes}

%------------------------------------------------------------------------------%
\begin{document}

\addtitle

%------------------------------------------------------------------------------%
\section{Diagonalização}

%------------------------------------------------------------------------------%
\begin{frame}{Matriz Diagonal}
  Uma matriz quadrada \(D\) é \emph{diagonal}

  \pause
  quando todos os elementos fora da diagonal principal são nulos
  \pause
  \[
    D \;=\;
    \begin{bmatrix}
      \lambda_{1}          & \z                   & {\color{gray}\cdots} & \z                   \\
      \z                   & \lambda_{2}          & {\color{gray}\cdots} & \z                   \\
      {\color{gray}\vdots} & {\color{gray}\vdots} & {\color{gray}\ddots} & {\color{gray}\vdots} \\
      \z                   & \z                   & {\color{gray}\cdots} & \lambda_{n}
    \end{bmatrix}
    \pause
    \;=\; \diag(\lambda_{1}, \lambda_{2}, \ldots, \lambda_{n})
  \]

  \bigskip
  \pause
  As operações com matrizes diagonais são especialmente simples
\end{frame}

%------------------------------------------------------------------------------%
\begin{frame}{Matrizes Semelhantes}
  Duas matrizes \(A\) e \(D\) são \emph{semelhantes}

  \pause
  quando existe uma matriz invertível \(P\)
  \pause
  \[
    D = P^{-1} A P
  \]
  \pause
  Equivalentemente
  \[
    A = P D P^{-1}
  \]
\end{frame}

%------------------------------------------------------------------------------%
\begin{frame}{Matrizes Diagonalizável}
  Uma matriz \(A\) é \emph{diagonalizável} quando é semelhante a uma matriz diagonal
\end{frame}

%------------------------------------------------------------------------------%
\begin{frame}{Autovetores e Diagonalização}
  Considere que $A_n$ possui $n$ \emph{autovetores} (\(A\vec{v}=\lambda\vec{v}\))

  \pause
  \emph{linearmente independentes}
  \[
    \vec{v}_{1},
    \vec{v}_{2},
    \ldots,
    \vec{v}_{n}
  \]
  \pause
  associados ao \emph{autovalores}
  \(
    \lambda_{1},
    \lambda_{2},
    \ldots,
    \lambda_{n}
  \)

  \pause
  Definimos
  \[
    P =
    \begin{bmatrix}[c]
      |              & |           &        & |              \\
      \; \vec{v}_{1} & \vec{v}_{2} & \cdots & \vec{v}_{n} \; \\
      |              & |           &        & |
    \end{bmatrix}
    \pause
    \qquad\text{ e }\qquad
    D = \diag(\lambda_{1}, \lambda_{2}, \ldots, \lambda_{n})
  \]
\end{frame}

%------------------------------------------------------------------------------%
\begin{proofframe}{Calculando $AP$}
  \begin{align*}
    \onslide<+->{AP}
    \onslide<+->{& = A
    \begin{bmatrix}[c]
      \; \vec{v}_{1} & \vec{v}_{2} & \cdots & \vec{v}_{n} \;
    \end{bmatrix}}
    \\
    \onslide<+->{& =
    \begin{bmatrix}[c]
      A \vec{v}_{1} & A \vec{v}_{2} & \cdots & A \vec{v}_{n}
    \end{bmatrix}}
    \\
    \onslide<+->{& =
    \begin{bmatrix}[c]
      \lambda_{1} \vec{v}_{1} & \lambda_{2} \vec{v}_{2} & \cdots & \lambda_{n} \vec{v}_{n}
    \end{bmatrix}}
    \\
    \onslide<+->{& = P \diag(\lambda_{1}, \lambda_{2}, \ldots, \lambda_{n})}
    \\[5mm]
    \onslide<+->{AP & = PD}
    \\
    \onslide<+->{A & = PDP^{-1}}
    \\
    \onslide<+->{P^{-1}AP & = D}
  \end{align*}
\end{proofframe}

%------------------------------------------------------------------------------%
\begin{frame}{Critério de Diagonalização}
\begin{enumerate}[<+->]
  \setlength{\itemsep}{2em}

  \item Uma matriz \(A\in\mathbb{R}^{n\times n}\) é diagonalizável
        \\ \medskip
        \onslide<+->{se, e somente se, possui \(n\) autovetores linearmente independentes}

  \item Se \(A\) possui \(n\) autovalores distintos, então é diagonalizável

  \item Autovalores repetidos não impedem a diagonalização
\end{enumerate}
\end{frame}

%------------------------------------------------------------------------------%
\section{Exemplos}

%------------------------------------------------------------------------------%
%------------------------------------------------------------------------------%
\begin{question}
  Diagonalize a matriz
  \[
    A=
    \begin{bmatrix}
      4 & 1 \\
      2 & 3
    \end{bmatrix}
  \]
  Determine \(P\) e \(D\) tais que
  \(
    A = PDP^{-1}
  \)
\end{question}

%------------------------------------------------------------------------------%
\begin{resolution}{Autovalores}
\begin{columns}
\column{0.4\textwidth}
  \onslide<+->{Polinômio característico}
  \begin{align*}
    \onslide<+->{p(\lambda)}
    \onslide<+->{& = \det(A-\lambda I)}
    \\
    \onslide<+->{& = \begin{vmatrix}[c]
      4-\lambda & 1         \\
      2         & 3-\lambda
    \end{vmatrix}}
    \\
  \onslide<+->{& = (4-\lambda)(3-\lambda) - 2 \\}
  \onslide<+->{& = 12 -7\lambda + \lambda^2 - 2 \\}
  \onslide<+->{& = \lambda^{2} - 7\lambda + 10}
  \end{align*}
\column{0.6\textwidth}
  \[
    \onslide<+->{\Delta}
    \onslide<+->{= b^2 - 4ac}
    \onslide<+->{= (-7)^2 - 4 \times 1 \times 10}
    \onslide<+->{= 9}
  \]

  \smallskip
  \[
    \onslide<+->{\sqrt{\Delta}}
    \onslide<+->{= \sqrt{9}}
    \onslide<+->{= 3}
  \]

  \smallskip
  \[
    \onslide<+->{\lambda_1}
    \onslide<+->{= \frac{-b+\sqrt{\Delta}}{2a}}
    \onslide<+->{= \frac{7+3}{2}}
    \onslide<+->{= 5}
  \]

  \smallskip
  \[
    \onslide<+->{\lambda_2}
    \onslide<+->{= \frac{-b-\sqrt{\Delta}}{2a}}
    \onslide<+->{= \frac{7-3}{2}}
    \onslide<+->{= 2}
  \]
\end{columns}
\end{resolution}

%------------------------------------------------------------------------------%
\begin{resolution}{Autovetores para $\lambda_1=5$}
\begin{columns}
\column{0.4\textwidth}
  \begin{align*}
    \onslide<+->{(A-5I)\vec{x}&=\vec{0}}
    \\
    \onslide<+->{\begin{bmatrix}[c]
      4 - 5& 1 \\
      2 & 3 - 5
    \end{bmatrix}
    \vetor{x}{y} & = \vetor{0}{0}}
    \\
    \onslide<+->{\begin{bmatrix}
      -1 &  1 \\
       2 & -2
    \end{bmatrix}
    \vetor{x}{y} & = \vetor{0}{0}}
  \end{align*}
  \begin{align*}
    \onslide<+->{-x + y & = 0} \\
    \onslide<+->{     y & = x}
  \end{align*}
\column{0.4\textwidth}
  \onslide<+->{Escolhendo $x=1$}
  \onslide<+->{\[
    y = 1
  \]}
  \onslide<+->{\[
    \vec{v}_1 = \vetor{1}{1}
  \]}
  \onslide<+->{\[
    E(5) = \ger \left\{ \vetor{1}{1} \right\}
  \]}
\end{columns}
\end{resolution}

%------------------------------------------------------------------------------%
\begin{resolution}{Autovetores para $\lambda_2=2$}
\begin{columns}
\column{0.4\textwidth}
  \begin{align*}
    \onslide<+->{(A-2I)\vec{x}&=\vec{0}}
    \\
    \onslide<+->{\begin{bmatrix}[c]
      4 - 2 & 1 \\
      2 & 3 - 2
    \end{bmatrix}
    \vetor{x}{y} & = \vetor{0}{0}}
    \\
    \onslide<+->{\begin{bmatrix}
      2 & 1 \\
      2 & 1
    \end{bmatrix}
    \vetor{x}{y} & = \vetor{0}{0}}
  \end{align*}
  \begin{align*}
    \onslide<+->{2x + y & = 0} \\
    \onslide<+->{     y & = -2x}
  \end{align*}
\column{0.4\textwidth}
  \onslide<+->{Escolhendo $x=1$}
  \onslide<+->{\[
    y = -2
  \]}
  \onslide<+->{\[
    \vec{v}_2 = \vetor{1}{-2}
  \]}
  \onslide<+->{\[
    E(2) = \ger \left\{ \vetor{1}{-2} \right\}
  \]}
\end{columns}
\end{resolution}

%------------------------------------------------------------------------------%
\begin{resolution}{Construção de \(P\) e \(D\)}
  Colocamos os autovetores nas colunas de \(P\)
  \pause
  \[
    P
    \pause
    = \left[ \vec{v}_2 \quad \vec{v}_2 \right]
    \pause
    = \begin{bmatrix}
      1 &  1  \\
      1 & -2
    \end{bmatrix}
  \]

  \pause
  Na mesma ordem dos autovalores de \(D\)
  \pause
  \[
    D
    \pause
    = \diag(\lambda_1, \lambda_2)
    \pause
    = \diag(5, 2)
    \pause
    = \begin{bmatrix}
        5 & 0 \\
        0 & 2
      \end{bmatrix}
  \]
\end{resolution}

%------------------------------------------------------------------------------%
\begin{proofframe}{Exemplo 1 -- Verificação -- Cálculo de \(P^{-1}\)}
  Para uma matriz de ordem \(2\)
  \[
    \begin{bmatrix}
      a & b \\
      c & d
    \end{bmatrix}^{-1}
    =
    \frac{1}{ad-bc}
    \begin{bmatrix}
      d  & -b \\
      -c & a
    \end{bmatrix}
  \]
  \pause
  Calculamos o determinante de \(P\)
  \pause
  \[
    \det(P)
    \pause
    = \begin{vmatrix}
      1 &  1  \\
      1 & -2
    \end{vmatrix}
    \pause
    = -2 -1
    \pause
    = -3
    \neq 0
  \]
  \pause
  \[
    P^{-1}
    \pause
    = \frac{1}{-3}
    \begin{bmatrix}
      -2 & -1 \\
      -1 & 1
    \end{bmatrix}
    \pause
    =
    \begin{bmatrix}
      \sfrac{2}{3} &  \sfrac{1}{3}  \\[3mm]
      \sfrac{1}{3} & -\sfrac{1}{3}
    \end{bmatrix}
  \]
\end{proofframe}

%------------------------------------------------------------------------------%
\begin{proofframe}{Exemplo 1 -- Verificação -- Diagonalização}
  \begin{align*}
    \onslide<+->{PDP^{-1}}
    \onslide<+->{& = \begin{bmatrix}
          1 & 1  \\
          1 & -2
        \end{bmatrix}
        \begin{bmatrix}
          5 & 0 \\
          0 & 2
        \end{bmatrix}
        P^{-1}}
    \\[2mm]
    \onslide<+->{& = \begin{bmatrix}
          5 & 2  \\
          5 & -4
        \end{bmatrix}
    P^{-1}}
    \\[2mm]
    \onslide<+->{& = \begin{bmatrix}
          5 & 2  \\
          5 & -4
        \end{bmatrix}
        \begin{bmatrix}
          \sfrac{2}{3} &  \sfrac{1}{3}  \\[3mm]
          \sfrac{1}{3} & -\sfrac{1}{3}
        \end{bmatrix}}
    \\[2mm]
    \onslide<+->{& = \begin{bmatrix}
          4 & 1 \\
          2 & 3
        \end{bmatrix}}
    \onslide<+->{= A}
  \end{align*}
\end{proofframe}

%------------------------------------------------------------------------------%
\endinput

%------------------------------------------------------------------------------%
\begin{question}
  Verifique se a matriz
  \[
    B=
    \begin{bmatrix}
      2 & 1 & 0 \\
      0 & 2 & 0 \\
      0 & 0 & 5
    \end{bmatrix}
  \]
  é diagonalizável
\end{question}

%------------------------------------------------------------------------------%
\begin{resolution}{Autovalores}
  Como a matriz é triangular
  \[
    p_{B}(\lambda)=(2-\lambda)^2(5-\lambda)
  \]
  Os autovalores são
  \[
    \lambda_{1}=2
    \qquad
    \lambda_{2}=5
  \]
  As multiplicidades algébricas são
  \[
    m_{a}(2)=2
    \qquad
    m_{a}(5)=1
  \]
\end{resolution}

%------------------------------------------------------------------------------%
\begin{resolution}{Autoespaços}
  Para \(\lambda=2\)
  \[
    (B-2I)\vec{v}=\vec{0}
  \]
  implica
  \[
    y=0
    \qquad
    z=0
  \]
  Logo
  \[
    E_{2}=\operatorname{span}
    \left\{
    \begin{bmatrix}
    1\\0\\0
    \end{bmatrix}
    \right\}
  \]
  Assim
  \[
    \dim(E_{2})=1
  \]
\end{resolution}

%------------------------------------------------------------------------------%
\begin{resolution}{Conclusão}
  Para \(\lambda=5\)
  \[
    E_{5}=\operatorname{span}
    \left\{
    \begin{bmatrix}
      0 \\
      0 \\
      1
    \end{bmatrix}
    \right\}
  \]
\end{resolution}

%------------------------------------------------------------------------------%
\begin{resolution}{Conclusão}
  Portanto, a soma das dimensões dos autoespaços é
  \[
    \dim(E_{2})+\dim(E_{5})=1+1=2
  \]
  Para diagonalizar uma matriz de ordem \(3\),

  precisamos de três autovetores linearmente independentes.

  Logo \(B\) não é diagonalizável
\end{resolution}

%------------------------------------------------------------------------------%
\endinput


%------------------------------------------------------------------------------%
\begin{question}
  Diagonalize, se possível, a matriz
  \[
    C=
    \begin{bmatrix}
      2 & 1 \\
      1 & 2
    \end{bmatrix}
  \]
\end{question}

%------------------------------------------------------------------------------%
\begin{resolution}{Autovalores}
  \begin{align*}
    \det(C-\lambda I)          & =0 \\
    (2-\lambda)^2-1            & =0 \\
    (2-\lambda-1)(2-\lambda+1) & =0 \\
    (1-\lambda)(3-\lambda)     & =0
  \end{align*}
  Logo
  \[
    \lambda_{1}=3
    \qquad
    \lambda_{2}=1
  \]
\end{resolution}

%------------------------------------------------------------------------------%
\begin{resolution}{Autovetores}
\begin{columns}
\column{0.4\textwidth}
  Para \(\lambda_{1}=3\)
  \begin{align*}
    (C-3I)\vec{v}_{1} & =\vec{0} \\
    -x+y              & =0       \\
    y                 & =x
  \end{align*}
  Escolhemos
  \[
    \vec{v}_{1}=
    \begin{bmatrix}
      1 \\
      1
    \end{bmatrix}
  \]
\column{0.4\textwidth}
  Para \(\lambda_{2}=1\)
  \begin{align*}
    (C-I)\vec{v}_{2} & =\vec{0} \\
    x+y              & =0       \\
    y                & =-x
  \end{align*}
  Escolhemos
  \[
    \vec{v}_{2}=
    \begin{bmatrix}
      1  \\
      -1
    \end{bmatrix}
  \]
\end{columns}
\end{resolution}

%------------------------------------------------------------------------------%
\begin{resolution}{Diagonalização}
  As matrizes são
  \[
    P=
    \begin{bmatrix}
      1 & 1  \\
      1 & -1
    \end{bmatrix}
    \qquad
    D=
    \begin{bmatrix}
      3 & 0 \\
      0 & 1
    \end{bmatrix}
  \]
\end{resolution}

%------------------------------------------------------------------------------%
\begin{resolution}{Diagonalização}
  Como
  \[
    \det(P)=-2\neq 0
  \]
  a matriz \(P\) é invertível

  Sua inversa é
  \[
    P^{-1}=
    \begin{bmatrix}
      \frac{1}{2} & \frac{1}{2}  \\
      \frac{1}{2} & -\frac{1}{2}
    \end{bmatrix}
  \]
  Portanto
  \[
    C=PDP^{-1}
  \]
\end{resolution}

%------------------------------------------------------------------------------%
\endinput


%%------------------------------------------------------------------------------%
%\section{Lista Mínima}
%
%%------------------------------------------------------------------------------%
%\begin{frame}{Lista Mínima}
%  \begin{enumerate}
%    \item
%  \end{enumerate}
%
%  \vfill
%  Atenção: A prova é baseada no livro, não nas apresentações
%\end{frame}

%------------------------------------------------------------------------------%
\end{document}
%------------------------------------------------------------------------------%
